% Target: Letters in Mathematical Physics
% Format: REVTeX4-2, two-column, PRL style, 4 pages strict
% Status: DRAFT (under review)
% Companion papers:
%   Paper 1: "Formal Verification of Alternative Mathematical Frameworks" (CICM 2026)
%   Paper 2: "Character Decomposition from Fortescue to Cartan-Weyl" (AACA, submitted)
%   Paper 3: "SO(14) Unification: A Machine-Verified Algebraic Scaffold" (PRD, in preparation)
%
% Claim tags:
%   [MV] Machine-verified in Lean 4 (0 sorry gaps)
%   [CO] Computed (Python scripts, analytical calculation)
%   [CP] Candidate physics (proposed, not proved)
%   [SP] Standard physics (textbook results)
%   [OP] Open problem (honestly stated)

\documentclass[aps,prl,twocolumn,showpacs,preprintnumbers,superscriptaddress,floatfix]{revtex4-2}

\usepackage{amsmath,amssymb}
\usepackage{mathtools}
\usepackage{booktabs}
\usepackage{hyperref}
\usepackage{xcolor}

% Notation
\newcommand{\so}{\mathfrak{so}}
\newcommand{\su}{\mathfrak{su}}
\newcommand{\Cl}{\mathrm{Cl}}
\newcommand{\Tr}{\mathrm{Tr}}
\newcommand{\ad}{\mathrm{ad}}
\newcommand{\MV}{\textsuperscript{[MV]}}
\newcommand{\CO}{\textsuperscript{[CO]}}
\newcommand{\CP}{\textsuperscript{[CP]}}
\newcommand{\SP}{\textsuperscript{[SP]}}
\newcommand{\OP}{\textsuperscript{[OP]}}

\begin{document}

\preprint{UFT-2026-04}

\title{Three Fermion Generations from $E_8$: Machine-Verified
Dimensional Constraints}

\author{Ian M.~Arnoldy}
\affiliation{Independent Researcher}

\date{March 10, 2026}

\begin{abstract}
We machine-verify the dimensional arithmetic underlying two results
about the three-generation problem in grand unified theories,
using the Lean~4 theorem prover~\cite{demoura2021lean4}
with the mathlib library~\cite{mathlib} across 9 proof files
(over 575 verified declarations, zero \texttt{sorry} gaps).
\emph{Obstruction}: assuming the standard spinor branching
rule\SP, the semi-spinor $\Delta^+(14)$ of $\mathrm{Spin}(14)$
decomposes under $\mathrm{SO}(10) \times \mathrm{SO}(4)$ as
$(\mathbf{16},\mathbf{2}) \oplus (\overline{\mathbf{16}},\mathbf{2})$,
giving multiplicity~2 for the $\mathbf{16}$.  The arithmetic
$3 \nmid 2$ is machine-verified\MV; the conclusion that no
$\mathrm{SO}(10{+}2k)$ spinor decomposition produces exactly three
generations follows from this arithmetic together with the
branching rule\SP.
\emph{Candidate resolution}\CP: $E_8$ admits an $\mathrm{SU}(9)/\mathbb{Z}_3$
decomposition where $248 = 80 + 84 + \overline{84}$\MV\ and
$\Lambda^3(\mathbb{C}^9) = 84$\MV.  The identification of
$3 \times 16 = 48$ of these dimensions with three Standard Model
generations requires the candidate-physics assumption that
$\mathrm{SU}(3)_{\mathrm{family}} \subset
\mathrm{SU}(4) \subset \mathrm{SU}(9)$ acts as a family
symmetry\CP.
The dimensional consistency of the two decompositions inside $E_8$
is verified: $49 + 21 + 21 = 91$\MV.
Every dimensional identity is verified by the Lean kernel;
every physical interpretation is explicitly tagged.
\end{abstract}

\pacs{12.10.Dm, 11.30.Hv, 11.15.Ex, 02.10.Ud}

\maketitle

% ============================================================================
\section{Introduction}
\label{sec:intro}
% ============================================================================

\emph{Scope note.}  The Lean~4 proofs in this Letter verify
\emph{dimensional arithmetic}: identities such as
$16 \times 2 + 16 \times 2 = 64$ and $\binom{9}{3} = 84$.
The representation-theoretic branching rules that give these
dimensions physical meaning are standard results\SP, cited from
the literature, not themselves machine-verified.  We distinguish
these levels throughout with explicit tags.

The Standard Model of particle physics contains three generations
of fermions---an empirical fact without theoretical explanation.
In grand unified theories (GUTs), the generation count depends on
the decomposition of matter representations under the unification
group.  The $\mathbf{16}$-dimensional chiral spinor of
$\mathrm{SO}(10)$~\cite{fritzsch1975,georgi1975} accommodates one
complete generation\SP; the question is how many copies appear.

For $\mathrm{SO}(14)$, which unifies gauge and gravitational degrees
of freedom via $\so(14) \supset \so(10) \oplus
\so(1,3)$~\cite{nesti2007,nesti2010,krasnov2018}, the semi-spinor
$\Delta^+(14)$ is 64-dimensional.  Under $\mathrm{SO}(10) \times
\mathrm{SO}(4)$, it decomposes with multiplicity~2 for the
$\mathbf{16}$~\cite{wilczek1982,slansky1981}.  The mismatch
$3 \nmid 2$ has been known since the 1980s~\cite{ida1980,ma1981,ramond1998}
but was never formally verified, and the resolution remained open.

We machine-verify the dimensional arithmetic in both results
using 9~Lean~4 proof files containing over 575 verified declarations
with zero \texttt{sorry} gaps:
\begin{enumerate}
  \item \emph{Obstruction}: assuming the standard spinor branching
    rule\SP, the family multiplicity $2^{k-1}$ is never~3\MV,
    so no $\mathrm{SO}(10{+}2k)$ gives three generations via
    spinor decomposition.
  \item \emph{Candidate resolution}\CP: the $\mathrm{SU}(9)/\mathbb{Z}_3$
    maximal subgroup of $E_8$ yields $\Lambda^3(\mathbb{C}^9)$
    with $\binom{9}{3} = 84$\MV, whose dimensional structure
    is consistent with three generations.
\end{enumerate}
To our knowledge, this is the first machine-verified dimensional
analysis of a three-generation mechanism in any GUT.  The
methodology---Lean~4 with mathlib---was developed for this
project~\cite{arnoldy2026cicm,arnoldy2026aaca,arnoldy2026prd}.

% ============================================================================
\section{Spinor Parity Obstruction}
\label{sec:obstruction}
% ============================================================================

The semi-spinor of $\mathrm{Spin}(2n)$ has dimension $2^{n-1}$.
For $\mathrm{Spin}(14)$ ($n = 7$): $\dim \Delta^+ = 2^6 = 64$\MV\
(\texttt{spinor\_branching\_64}).

Under the maximal subgroup $\mathrm{SO}(10) \times \mathrm{SO}(4)
\subset \mathrm{SO}(14)$, the standard spinor branching
rule~\cite{slansky1981} gives
\begin{equation}
  \Delta^+(14) = (\mathbf{16}, \mathbf{2})
  \oplus (\overline{\mathbf{16}}, \mathbf{2}),
  \label{eq:branching}
\end{equation}
where $\mathbf{16} = \Delta^+(10)$ (one SM generation\SP) and
$\mathbf{2}$ is a spinor of $\mathrm{SO}(4)$.  The dimension
check $16 \times 2 + 16 \times 2 = 64$ is machine-verified\MV\
(\texttt{branching\_rule\_dimension}).

The multiplicity of the $\mathbf{16}$ equals the dimension of
the $\mathrm{SO}(4)$ spinor:
\begin{equation}
  \mathrm{mult}(\mathbf{16}) = \dim \Delta^+(\mathrm{SO}(4))
  = 2^{2-1} = 2.
  \label{eq:mult}
\end{equation}
This is verified as \texttt{multiplicity\_is\_two}\MV.

The obstruction is arithmetic: $3 \nmid 2$\MV\
(\texttt{three\_gen\_excluded}).  Combined with the branching
rule~\eqref{eq:branching}\SP, this excludes exactly three copies
of the $\mathbf{16}$ from $\Delta^+(14)$ under any
$\mathrm{SO}(10) \times \mathrm{SO}(4)$ decomposition.

Two extensions strengthen the result:

\emph{Universality.}  All embeddings $\mathrm{SO}(10) \times
\mathrm{SO}(4) \hookrightarrow \mathrm{SO}(14)$ are conjugate
under $\mathrm{SO}(14)$\SP, because $\mathrm{SO}(14)$ acts
transitively on the Grassmannian $\mathrm{Gr}(10,14)$.
The orbit-stabilizer dimension $\dim \mathrm{Gr}(10,14) =
91 - 51 = 40$ is arithmetically verified\MV\
(\texttt{orbit\_stabilizer\_check}); the conjugacy itself is
standard Lie theory\SP.
No non-standard embedding escapes the obstruction.

\emph{No orthogonal escape.}  For any $\mathrm{SO}(10{+}2k)$,
the family multiplicity is $2^{k-1}$.  We prove that $2^k \neq 3$
for all $k \in \mathbb{N}$\MV\ (\texttt{no\_so\_gives\_three\_families}),
since 3~is not a power of~2.  The Dirac (128-dimensional) spinor
gives multiplicity~4, and $3 \nmid 4$\MV\
(\texttt{three\_gen\_excluded\_dirac}).

The capstone theorem \texttt{spinor\_parity\_obstruction}\MV\
collects all four ingredients---dimension, multiplicity, divisibility,
and dimension consistency---into a single conjunction:
\begin{equation}
  \Delta^+(14) = 64,\;
  \mathrm{mult} = 2,\;
  3 \nmid 2,\;
  16 \cdot 2 + 16 \cdot 2 = 64.
  \label{eq:capstone-obs}
\end{equation}

Given the standard branching rule\SP, the dimensional arithmetic
excludes three fermion generations from any
$\mathrm{SO}(10{+}2k)$ spinor decomposition\MV.  If the
generation mechanism exists, it must be extrinsic to the
orthogonal series.

% ============================================================================
\section{Three Generations from $E_8$}
\label{sec:e8}
% ============================================================================

$E_8$ provides the extrinsic structure.  The chain proceeds in
three steps.

\subsection{Embedding: $\mathrm{SO}(14) \subset \mathrm{SO}(16) \subset E_8$}

$E_8$ has dimension 248 and contains $\mathrm{SO}(16)$ as a
maximal subgroup\SP.  Under $\mathrm{SO}(16)$, the adjoint
decomposes as~\cite{slansky1981}
\begin{equation}
  248 = 120 + 128
  = \ad(\mathrm{SO}(16)) \oplus S^+(\mathrm{Spin}(16)),
  \label{eq:e8-so16}
\end{equation}
verified as \texttt{e8\_so16\_decomposition}\MV.
$\mathrm{SO}(14) \times \mathrm{SO}(2)$ sits block-diagonally
inside $\mathrm{SO}(16)$: $120 = 91 + 1 + 28$\MV.
The semi-spinor branches as $128 = 64 + 64$\MV.  Together:
\begin{equation}
  248 = 91 + 1 + 28 + 64 + 64,
  \label{eq:e8-so14}
\end{equation}
verified as
\texttt{e8\_so14\_full\_decomposition}\MV, with the full chain
collected in \texttt{e8\_embedding\_chain\_complete}\MV.

\subsection{$\mathrm{SU}(9)/\mathbb{Z}_3$ decomposition}

$E_8$ admits a \emph{second} maximal subgroup decomposition under
$\mathrm{SU}(9)/\mathbb{Z}_3$~\cite{slansky1981,adams1996,ramond2010}:
\begin{equation}
  248 = 80 + 84 + \overline{84}
  = \ad(\mathrm{SU}(9)) \oplus \Lambda^3(\mathbb{C}^9) \oplus
  \overline{\Lambda^3(\mathbb{C}^9)},
  \label{eq:e8-su9}
\end{equation}
where $80 = 9^2 - 1$ and $84 = \binom{9}{3}$.  Both summands are
verified: \texttt{e8\_su9\_decomposition}\MV.

The 84 and $\overline{84}$ are conjugate under Hodge duality:
$\binom{9}{3} = \binom{9}{6}$\MV.  The $\mathbb{Z}_3$
quotient arises because the center $\mathbb{Z}_9$ of
$\mathrm{SU}(9)$ acts by cube roots of unity on the three
eigenspaces (80, 84, $\overline{84}$)\SP.

\subsection{Cauchy decomposition and the family triplet}

The fundamental $\mathbf{9}$ of $\mathrm{SU}(9)$ branches under
$\mathrm{SU}(5) \times \mathrm{SU}(4) \subset \mathrm{SU}(9)$~\cite{georgi1974,slansky1981} as
$\mathbf{9} = (\mathbf{5},\mathbf{1}) \oplus (\mathbf{1},\mathbf{4})$\SP.
The Cauchy formula for exterior powers
$\Lambda^k(V \oplus W) = \bigoplus_{i+j=k} \Lambda^i(V) \otimes
\Lambda^j(W)$ gives:
\begin{equation}
  84 = \underbrace{\binom{5}{3}}_{10} +
  \underbrace{\binom{5}{2}\binom{4}{1}}_{40} +
  \underbrace{\binom{5}{1}\binom{4}{2}}_{30} +
  \underbrace{\binom{4}{3}}_{4},
  \label{eq:cauchy}
\end{equation}
verified as \texttt{cauchy\_decomp\_84}\MV.  In representation
notation: $84 = (\mathbf{10},\mathbf{1}) \oplus
(\mathbf{10},\mathbf{4}) \oplus (\mathbf{5},\mathbf{6}) \oplus
(\mathbf{1},\overline{\mathbf{4}})$.

Now break $\mathrm{SU}(4) \to \mathrm{SU}(3)_\mathrm{family}
\times \mathrm{U}(1)$\CP, under which $\mathbf{4} = \mathbf{3}
+ \mathbf{1}$.  The $(\mathbf{10},\mathbf{4})$ piece branches as
$(\mathbf{10},\mathbf{3}) \oplus (\mathbf{10},\mathbf{1})$, giving
\emph{three copies} of the $\mathbf{10}$ of $\mathrm{SU}(5)$:
$\dim = 10 \times 3 = 30$\MV\ (\texttt{three\_generation\_count}).

\begin{table}[t]
\caption{Generation content from $84 \oplus \overline{84}$ under
$\mathrm{SU}(5) \times \mathrm{SU}(3)_\mathrm{family}$.
The three generations ($3 \times 16 = 48$ dimensions) are drawn from
\emph{both} $\mathbb{Z}_3$ sectors: the $\mathbf{10}$ and
$\mathbf{1}$ from the 84 ($J = +i$ eigenspace), the
$\overline{\mathbf{5}}$ from the $\overline{84}$ ($J = -i$).
Remaining dimensions ($84 + 84 - 48 = 120$) are exotic.
All dimensions are machine-verified\MV.}
\label{tab:84}
\begin{ruledtabular}
\begin{tabular}{lcccc}
Rep.\ $(\mathrm{SU}(5), \mathrm{SU}(3)_f)$ & Dim & Source & $J$ & Content \\
\colrule
$(\mathbf{10},\mathbf{3})$ & 30 & 84 & $+i$ & $3\times\mathbf{10}$ \\
$(\mathbf{1},\bar{\mathbf{3}})$ & 3  & 84 & $+i$ & $3\times\nu^c$ \\
$(\bar{\mathbf{5}},\bar{\mathbf{3}})$ & 15 & $\overline{84}$ & $-i$ & $3\times\bar{\mathbf{5}}$ \\
\colrule
Generation total & 48 & & & $3 \times 16$ \\
Exotic remainder & 120 & & & vector-like pairs \\
\end{tabular}
\end{ruledtabular}
\end{table}

Similarly, $(\mathbf{5},\mathbf{6}) \to (\mathbf{5},\mathbf{3})
\oplus (\mathbf{5},\overline{\mathbf{3}})$ and
$(\mathbf{1},\overline{\mathbf{4}}) \to
(\mathbf{1},\overline{\mathbf{3}}) \oplus (\mathbf{1},\mathbf{1})$.
The conjugate $\overline{84}$ decomposes analogously;
$(\overline{\mathbf{5}},\overline{\mathbf{3}})$ with dimension~15
appears in the $\overline{84}$, providing $3 \times \overline{\mathbf{5}}$.
Table~\ref{tab:84} collects the generation content from both sectors.
Three complete SM generations contribute
$3 \times (\mathbf{10} + \overline{\mathbf{5}} + \mathbf{1}) =
3 \times 16 = 48$ dimensions\MV\
(\texttt{three\_gen\_dim}, \texttt{generation\_matter\_total}),
drawn from both the 84 and $\overline{84}$.
The exotic content is $84 + 84 - 48 = 120$ dimensions,
forming vector-like pairs that can acquire GUT-scale masses\CP.

The eight-part capstone \texttt{e8\_three\_generation\_theorem}\MV\
collects: $\binom{9}{3} = 84$, $10 \times 3 = 30$,
$5 \times 3 = 15$, $10 + 5 + 1 = 16$, $3 \times 16 = 48$,
$84 - 48 = 36$, and $48 + 36 = 84$.
The anomaly-free eigenspace theorem
\texttt{j\_sector\_anomaly\_free}\MV\ additionally verifies
$\Tr(Y) = 0$ for each $J$ eigenspace independently.

\subsection{Compatibility of the two decompositions}

The $\mathrm{SO}(16)$ and $\mathrm{SU}(9)$ decompositions of $E_8$ share a common
subgroup $\mathrm{SU}(7) \times \mathrm{U}(1)$ (dimension
$48 + 1 = 49$\MV).  The 91 generators of $\mathrm{SO}(14)$
distribute across the $\mathrm{SU}(9)$ sectors as:
\begin{equation}
  91 = \underbrace{49}_{\text{in } 80} +
  \underbrace{21}_{\text{in } 84} +
  \underbrace{21}_{\text{in } \overline{84}},
  \label{eq:bridge}
\end{equation}
where $21 = \binom{7}{2} = \dim \Lambda^2(\mathbb{C}^7)$\MV.
This is verified as \texttt{so14\_compatible\_split}\MV.
Every generator of $\mathrm{SO}(14)$ is accounted for in
both decompositions.

The complete dimensional chain is collected in the seven-part
\texttt{three\_generation\_theorem}\MV, which verifies the arithmetic:
(a)~$3 \nmid 2$ (obstruction at $\mathrm{SO}(14)$),
(b)~$91 + 1 + 28 + 64 + 64 = 248$ (embedding dimensions),
(c)~$80 + 84 + 84 = 248$ ($\mathrm{SU}(9)$ decomposition),
(d)~$3 \times 16 = 48 \leq 84$ (generation count),
(e)~$84 - 48 = 36$ (exotics),
(f)~$49 + 21 + 21 = 91$ (compatibility), and
(g)~$3 \mid 3$ and $3 \nmid 2$ (origin of~3).

% ============================================================================
\section{Discussion}
\label{sec:discussion}
% ============================================================================

\emph{Epistemological boundary.}
Every dimensional identity in this Letter is machine-verified
(\MV): over 575~verified declarations across 9~Lean~4 files, zero \texttt{sorry}
gaps.  The arithmetic $\binom{9}{3} = 84$, the branching
$84 = 10 + 40 + 30 + 4$, and the generation count $10 \times 3 = 30$
are mathematical facts accepted by the Lean kernel.

The identification of $\mathrm{SU}(3)_\mathrm{family}$ as the
\emph{physical} family symmetry is candidate physics\CP.  The claim
that exotics decouple at the GUT scale is candidate physics\CP.
All Yukawa textures and proton decay predictions from the
$E_8$ breaking chain are open problems\OP.  We tag every claim
explicitly.

\emph{Chirality boundary.}
The dimensional analysis above is algebraic: it holds for
\emph{all} real forms of $E_8$, including the physically relevant
$E_8(-24)$ containing $\mathrm{Spin}(3{,}11)$.  The Clifford periodicity
distinguishes signatures---$\Cl(14,0)$ has $p{-}q \bmod 8 = 6$ while
$\Cl(11,3)$ has $p{-}q \bmod 8 = 0$---but the generation count
$3 \times 16 = 48$ is signature-independent\MV.
Distler and Garibaldi~\cite{distler2010} proved that no embedding of
the Standard Model in any real form of $E_8$ can produce three
\emph{massless chiral} generations from the adjoint.
Wilson~\cite{wilson2024b} argues that a compact gauge group is
intrinsically massless, so for theories with explicitly non-zero
masses---such as $E_8(-24)$, which contains $\mathrm{SU}(2,2)$ rather
than the Poincar\'e group---the Distler--Garibaldi no-go theorem does
not apply. The non-compact degrees of freedom serve as mass parameters
rather than ghosts.
Which definition of chirality is physically correct remains open\OP,
but the \emph{algebraic} question can be resolved.

We provide a two-level algebraic definition.
\emph{Level~1} (Definition~B): $\Lambda^3(\mathbb{C}^9)$ is \emph{not}
self-conjugate as an $\mathrm{SU}(9)$ representation, since
self-conjugacy of $\Lambda^k(\mathbb{C}^n)$ requires $2k = n$, and
$2 \times 3 = 6 \neq 9$\MV\ (\texttt{wilson\_chirality}).
Wilson~\cite{wilson2024b} makes this argument explicitly: ``Since
$\Lambda^3(\mathbb{C}^9)$ is not the same as the antisymmetric cube of
$9^*$, the theory is chiral.''
Because 9 is odd, \emph{every} exterior power $\Lambda^k(\mathbb{C}^9)$
for $0 < k < 9$ is non-self-conjugate\MV\
(\texttt{odd\_nine\_all\_chiral}); the chirality of $\mathrm{SU}(9)$
is structurally robust.
\emph{Level~2} (Definition~D): the $A_8$/$D_8$ overlap
matrix---computed from the non-nested maximal-rank subalgebras
$\mathrm{SU}(9)$ and $\mathrm{SO}(16)$ of $E_8$---quantifies
the asymmetry\CO: the 84 roots of sector~1 split as $28 + 35 + 21$
under $D_8 \supset D_7$ chirality (S$^+$ vs.\ S$^-$)\CO, giving a
chirality index $\chi = 35 - 21 = 14$\MV\
(\texttt{chirality\_index\_sector1})---the subtraction is
machine-verified, while the root classification is computed\CO.
The total 248 remains balanced ($64^+ = 64^-$)\MV, consistent
with Distler--Garibaldi.
The combined definition is collected in
\texttt{massive\_chirality\_definition}\MV\ (12-part conjunction)
and the full certificate \texttt{complete\_massive\_chirality\_skeleton}\MV\
(20-part conjunction), documented in
\texttt{massive\_chirality\_definition.lean},
\texttt{e8\_chirality\_boundary.lean}, and
\texttt{exterior\_cube\_chirality.lean}.

\emph{Anomaly cancellation and the $\mathbb{Z}_3$ sector operator.}
The $\mathrm{SU}(5)$ anomaly cancels generation by generation~\cite{georgi1982,baez2010}:
$A(\mathbf{10}) + A(\overline{\mathbf{5}}) = 1 + (-1) = 0$\SP.
For three generations: $3 \times 0 = 0$\MV\
(\texttt{anomaly\_three\_generations}).
The $\mathbb{Z}_3$ automorphism $\sigma$ of $E_8$ that defines the
$\mathrm{SU}(9)/\mathbb{Z}_3$ decomposition induces a complex structure
$J = (2/\sqrt{3})(\sigma + \tfrac{1}{2}\mathrm{Id})|_V$ on the
168-dimensional matter space $V = 84 \oplus \overline{84}$, with
$J^2 = -\mathrm{Id}$\CO.
$J$ does \emph{not} reduce to the chirality operator $\gamma_5$\CO:
one Standard Model generation straddles both $J$ eigenspaces, with the
$\mathbf{10} + \mathbf{1}$ (containing $Q_L, u^c, e^c, \nu^c$) in the
$+i$ eigenspace and the $\overline{\mathbf{5}}$ (containing $d^c, L$) in
the $-i$ eigenspace.
Instead, $J$ encodes the $\mathrm{SU}(5)$ multiplet structure:
it is the algebraic reason why fermions are grouped into the
$\mathbf{10}$ and $\overline{\mathbf{5}}$.
Each $J$ eigenspace is independently anomaly-free:
$\Tr(Y) = 0$ for $\{\mathbf{10},\mathbf{1}\}$ and for
$\{\overline{\mathbf{5}}\}$ separately\MV\
(\texttt{j\_sector\_anomaly\_free}).
This cancellation---previously an unexplained
feature of the Standard Model---becomes, in this framework,
a consequence of $\sigma$ preserving the algebraic structure
of $E_8$\CP.

\emph{Why $E_8$ succeeds where $\mathrm{SO}(14)$ fails.}
The root cause is algebraic: spinor representations of orthogonal
groups yield multiplicities $2^{k-1}$, which are always powers
of~2.  The exterior algebra $\Lambda^k(\mathbb{C}^N)$ has no such
constraint---$\binom{N}{k}$ can be any integer.
In particular, $\binom{3}{1} = 3$\MV, while no $2^k = 3$\MV\
(\texttt{three\_not\_power\_of\_two}).
The transition from $\mathrm{SO}(14)$ to $E_8$ replaces the $\mathrm{SO}(4)$ family
factor (spinor dim~2) with $\mathrm{SU}(3)_\text{family}$ (fundamental
dim~3).  The extra structure has dimension $\dim \mathrm{SU}(4) -
\dim \mathrm{SO}(4) = 15 - 6 = 9$\MV, which is dimensionally
consistent with $\mathrm{SU}(3) \times \mathrm{U}(1)$
($8 + 1 = 9$\MV); the identification as a physical family
symmetry is candidate physics\CP.

\emph{Comparison with prior work.}
Wilson~\cite{wilson2024,wilson2025} identified discrete order-3
elements of $E_8$ whose centralizer is $\mathrm{SU}(9)/\mathbb{Z}_3$,
providing a fundamentally discrete three-generation mechanism.
Our contribution is to \emph{machine-verify} the complete dimensional
chain from obstruction through candidate resolution.
Our chirality boundary analysis, informed by
Distler and Garibaldi~\cite{distler2010}, clarifies the precise
scope of Wilson's escape: the algebraic prerequisites are verified,
while the chirality definition remains open\OP.
Calabi-Yau compactification~\cite{candelas1985} obtains three
generations from the Euler characteristic $|\chi|/2 = 3$ of specific
manifolds, but this is geometric, not algebraic.
BenTov and Zee~\cite{bentov2016} explored $\mathrm{SO}(18)$ and showed
that the family problem persists for all $\mathrm{SO}(10{+}2k)$; our
\texttt{no\_so\_gives\_three\_families}\MV\ proves this universally.
Bars and G\"unaydin~\cite{bars1980} were the first to use $E_8$
for grand unification; our work adds formal verification.

\emph{Exotics.}
The 120 exotic dimensions in $84 \oplus \overline{84}$
(Table~\ref{tab:84}) form vector-like pairs under $\mathrm{SU}(5)$
that can acquire GUT-scale masses without breaking the Standard
Model gauge symmetry\CP.
The ratio of SM matter to total ($48/168 \approx 29\%$) is
typical of $E_8$ models~\cite{lisi2007}.

% ============================================================================
\section{Conclusion}
\label{sec:conclusion}
% ============================================================================

We have machine-verified the dimensional arithmetic showing that
three fermion generations are excluded from any
$\mathrm{SO}(10{+}2k)$ spinor decomposition (given the standard
branching rule\SP), and that the $E_8$ decomposition under
$\mathrm{SU}(9)/\mathbb{Z}_3$ is dimensionally consistent with
three generations ($3 \times 16 = 48 \leq 84$\MV).
Whether this dimensional consistency constitutes a physical
resolution depends on the candidate-physics identification of
$\mathrm{SU}(3)_\mathrm{family}$\CP.
We further provide a two-level algebraic definition of chirality:
non-self-conjugacy of $\Lambda^3(\mathbb{C}^9)$ (Definition~B\MV),
quantified by a chirality index $\chi = 14$ per $\mathbb{Z}_3$
sector (Definition~D; root classification\CO, subtraction\MV),
compatible with the Distler--Garibaldi no-go theorem.
The physical observable corresponding to massive
chirality remains open\OP.

The number~3 enters through the breaking chain
$E_8 \supset \mathrm{SU}(9) \supset \mathrm{SU}(5) \times
\mathrm{SU}(4) \supset \mathrm{SU}(5) \times
\mathrm{SU}(3)_\mathrm{family} \times \mathrm{U}(1)$\CP,
where 3 is the dimension of the fundamental of
$\mathrm{SU}(3)_\mathrm{family}$.
The Lean proof verifies the dimensional identity
$\dim(\mathrm{fund.\ SU}(3)) = 3$ by definition;
the claim that this $\mathrm{SU}(3)$ \emph{is} the family
symmetry is candidate physics\CP, not a verified consequence.

The complete argument---9 proof files, over 575 verified declarations, zero
\texttt{sorry} gaps---is available from the corresponding
author upon reasonable request~\cite{uft-github}.
We propose that machine verification become standard methodology
for GUT algebraic constructions, eliminating the class of errors
that can persist undetected for decades.

\begin{acknowledgments}
The machine verification was performed using Lean~4 and the mathlib
library.  The author thanks the Lean and mathlib communities for
providing the infrastructure that made this work possible.
Proof engineering and manuscript preparation were assisted by Claude
(Anthropic); all mathematical content, research direction, and
physical interpretation are the sole responsibility of the author.
\end{acknowledgments}

\section*{Data Availability}
The Lean~4 source code for all machine-verified proofs is available
from the corresponding author upon reasonable request~\cite{uft-github}.

% ============================================================================
% References
% ============================================================================

\begin{thebibliography}{29}

% === Grand unification ===

\bibitem{fritzsch1975}
H.~Fritzsch and P.~Minkowski,
``Unified interactions of leptons and hadrons,''
Ann.\ Phys.\ \textbf{93}, 193 (1975).

\bibitem{georgi1975}
H.~Georgi,
``The state of the art --- gauge theories,''
AIP Conf.\ Proc.\ \textbf{23}, 575 (1975).

\bibitem{georgi1974}
H.~Georgi and S.~L.~Glashow,
``Unity of all elementary-particle forces,''
Phys.\ Rev.\ Lett.\ \textbf{32}, 438 (1974).

\bibitem{georgi1982}
H.~Georgi,
\textit{Lie Algebras in Particle Physics}
(Benjamin/Cummings, 1982).

% === SO(14) / Gravity-gauge ===

\bibitem{nesti2007}
F.~Nesti and R.~Percacci,
``Graviweak unification,''
J.\ Phys.\ A \textbf{41}, 075405 (2008)
[arXiv:0706.3307].

\bibitem{nesti2010}
F.~Nesti and R.~Percacci,
``Chirality in unified theories of gravity,''
Phys.\ Rev.\ D \textbf{81}, 025010 (2010)
[arXiv:0909.4537].

\bibitem{krasnov2018}
K.~Krasnov and R.~Percacci,
``Gravity and unification: a review,''
Class.\ Quantum Grav.\ \textbf{35}, 143001 (2018)
[arXiv:1712.03061].

% === Historical three-generation ===

\bibitem{wilczek1982}
F.~Wilczek and A.~Zee,
``Families from spinors,''
Phys.\ Rev.\ D \textbf{25}, 553 (1982).

\bibitem{ida1980}
M.~Ida, Y.~Kayama, and T.~Kitazoe,
``Inclusion of generations in $\mathrm{SO}(14)$,''
Prog.\ Theor.\ Phys.\ \textbf{64}, 1745 (1980).

\bibitem{ma1981}
Z.-Q.~Ma, D.-S.~Du, P.-Y.~Xue, and X.-J.~Zhou,
``Generation problem and $\mathrm{SO}(14)$ grand unified theories,''
Chinese Physics C \textbf{5}, 664 (1981).

\bibitem{ramond1998}
P.~Ramond,
``The family problem,''
hep-ph/9809459 (1998).

% === E_8 ===

\bibitem{bars1980}
I.~Bars and M.~G\"unaydin,
``Grand unification with the exceptional group $E_8$,''
Phys.\ Rev.\ Lett.\ \textbf{45}, 859 (1980).

\bibitem{wilson2024}
R.~A.~Wilson,
``Uniqueness of an $E_8$ model of elementary particles,''
J.\ Math.\ Phys.\ \textbf{63}, 081703 (2022).

\bibitem{wilson2025}
R.~A.~Wilson,
``A synthesis of $E_8$ models of elementary particles,''
arXiv:2507.16517 (2025).

\bibitem{lisi2007}
A.~G.~Lisi,
``An exceptionally simple theory of everything,''
arXiv:0711.0770 (2007).

\bibitem{adams1996}
J.~F.~Adams,
\textit{Lectures on Exceptional Lie Groups}
(University of Chicago Press, 1996).

% === Three-generation mechanisms ===

\bibitem{bentov2016}
Y.~BenTov and A.~Zee,
``Origin of families and $\mathrm{SO}(18)$ grand unification,''
Phys.\ Rev.\ D \textbf{93}, 065036 (2016)
[arXiv:1505.04312].

\bibitem{candelas1985}
P.~Candelas, G.~T.~Horowitz, A.~Strominger, and E.~Witten,
``Vacuum configurations for superstrings,''
Nucl.\ Phys.\ B \textbf{258}, 46 (1985).

% === Group theory references ===

\bibitem{slansky1981}
R.~Slansky,
``Group theory for unified model building,''
Phys.\ Rep.\ \textbf{79}, 1 (1981).

\bibitem{ramond2010}
P.~Ramond,
\textit{Group Theory: A Physicist's Survey}
(Cambridge University Press, 2010).

\bibitem{baez2010}
J.~C.~Baez and J.~Huerta,
``The algebra of grand unified theories,''
Bull.\ Amer.\ Math.\ Soc.\ \textbf{47}, 483 (2010).

\bibitem{distler2010}
J.~Distler and S.~Garibaldi,
``There is no `Theory of Everything' inside $E_8$,''
Comm.\ Math.\ Phys.\ \textbf{298}, 419 (2010)
[arXiv:0904.1447].

\bibitem{wilson2024b}
R.~A.~Wilson,
``On the problem of interacting particles in theoretical physics,''
arXiv:2407.18279 (2024).

% === Formal verification ===

\bibitem{demoura2021lean4}
L.~de Moura and S.~Ullrich,
``The Lean~4 theorem prover and programming language,''
CADE-28, LNCS \textbf{12699}, 625 (2021).

\bibitem{mathlib}
The mathlib community,
``The Lean mathematical library,''
CPP 2020 (2020).

% === Companion papers ===

\bibitem{arnoldy2026cicm}
I.~M.~Arnoldy,
``Formal verification of alternative mathematical frameworks:
a case study using Lean~4,''
submitted to CICM 2026 (2026).

\bibitem{arnoldy2026aaca}
I.~M.~Arnoldy,
``Character decomposition from Fortescue to Cartan--Weyl:
unifying symmetrical components and root space decomposition
via Clifford algebras,''
submitted to Advances in Applied Clifford Algebras (2026).

\bibitem{arnoldy2026prd}
I.~M.~Arnoldy,
``SO(14) unification: a machine-verified algebraic scaffold
with phenomenological predictions,''
in preparation (2026).

\bibitem{uft-github}
I.~M.~Arnoldy,
``UFT: Machine-verified algebraic scaffold,''
\url{https://github.com/iarnoldy/UFT} (2026).

\end{thebibliography}

\end{document}
